\section{Conclusion}

This study examined how accessibility, motility, and time are represented in Swiss spatial planning discourse using a textometric approach. The analysis reveals a clear imbalance in planning language. Accessibility is widely present and discussed in concrete terms, whereas motility, which encompasses skills, motivation, and cognitive appropriation, is almost entirely absent. This absence indicates that capability-based dimensions of mobility are not yet part of the institutional vocabulary, making them difficult to plan for or evaluate.

Time is represented asymmetrically as well. References to physical and objective time appear frequently and are treated as operational variables to be measured or minimized. Perceived or subjective time appears to a lesser extent. Concepts such as comfort are mentioned but remain vague and underdeveloped. This suggests that planners have begun to acknowledge experiential aspects of travel, yet they do not articulate them in operational terms.

Methodologically, the textometric pipeline proved effective at revealing patterns of presence and absence in institutional language. It is valuable for mapping what planning discourse renders thinkable and actionable. At the same time, the approach is essentially descriptive. Combining it with contextual and semantic analyses would better recover meaning beyond raw counts, reducing reliance on manual interpretation and mitigating researcher bias.

The project’s contributions are twofold. First, it documents a persistent discursive bias in official planning documents toward supply-side and technical framings of mobility. Second, it provides a reproducible dataset and analysis pipeline, released under a CC-BY license in the project repository \cite{morand_nathmoswissspatialplanningtextometricanalysis_2025}, enabling further research or replication. Making this corpus freely available lowers the "cost" of entry and allows exploration of broader aspects of spatial planning beyond mobility alone.

Looking forward, bridging the gap between accessibility and motility requires both conceptual and institutional change. We need ways to measure motility without oversimplifying it, and planning should routinely collect data on experiential and behavioral dimensions without falling into's Goodhart's law. Only then can mobility policy move beyond networks to truly support people’s capabilities and maybe let the author sleep a little easier knowing that Switzerland point a bit more toward a sustainable future.